\documentclass[tikz,border=8pt]{standalone}
\usepackage{tikz}
\usepackage{amsmath}
\usepackage{amssymb}
\usepackage{pgfplots}
\pgfplotsset{compat=1.18}
\usepackage{ctex}
\usetikzlibrary{calc,positioning,arrows.meta,matrix,trees}

% LC 50 快速幂
% 计算 x^13，13 的二进制 = 1101
% 分治: x^13 = x^8 * x^4 * x^1（对应 bit3,bit2,bit0 为 1）

\begin{document}
\begin{tikzpicture}[
  every node/.style={font=\small},
  bitcell/.style={draw=gray!50, rectangle, minimum width=0.55cm, minimum height=0.55cm, thin, fill=white},
  onebit/.style={draw=blue!60, rectangle, minimum width=0.55cm, minimum height=0.55cm, thin, fill=blue!15},
  stepbox/.style={draw=blue!50, rectangle, rounded corners=2pt, minimum width=3.0cm, minimum height=0.65cm, thick, fill=blue!8},
  activebox/.style={draw=orange!70, rectangle, rounded corners=2pt, minimum width=3.0cm, minimum height=0.65cm, thick, fill=orange!10},
  resbox/.style={draw=green!60!black, rectangle, rounded corners=2pt, minimum width=3.0cm, minimum height=0.65cm, thick, fill=green!10},
  ptr/.style={->,>=Stealth,thick},
]

%% 标题
\node[font=\large\bfseries] at (5.5, 1.5)
  {快速幂（LC 50 Pow(x,n)）};
\node[font=\small,gray] at (5.5, 0.9)
  {计算 $x^{13}$，二进制 13 = 1101，分治跳过 bit=0 的步骤};

%% 二进制展示（13 = 1101）
\node[font=\scriptsize,gray] at (2.8, 0.2) {13 二进制：};
% bit3=1, bit2=1, bit1=0, bit0=1 -> 1,1,0,1 positions: i=0,1,2,3
\node[draw=blue!60,fill=blue!15,rectangle,minimum width=0.55cm,minimum height=0.55cm,thin]
  at (4.2+0*0.65, 0.2) {1};
\node[draw=blue!60,fill=blue!15,rectangle,minimum width=0.55cm,minimum height=0.55cm,thin]
  at (4.2+1*0.65, 0.2) {1};
\node[draw=gray!50,fill=white,rectangle,minimum width=0.55cm,minimum height=0.55cm,thin]
  at (4.2+2*0.65, 0.2) {0};
\node[draw=blue!60,fill=blue!15,rectangle,minimum width=0.55cm,minimum height=0.55cm,thin]
  at (4.2+3*0.65, 0.2) {1};
\node[font=\tiny,gray] at (4.2+0*0.65, -0.2) {$2^3$};
\node[font=\tiny,gray] at (4.2+1*0.65, -0.2) {$2^2$};
\node[font=\tiny,gray] at (4.2+2*0.65, -0.2) {$2^1$};
\node[font=\tiny,gray] at (4.2+3*0.65, -0.2) {$2^0$};
\node[font=\scriptsize,gray,anchor=west] at (6.95, 0.2) {= 8+4+1 = 13};

%% 快速幂迭代过程（表格，5步）
\node[font=\small\bfseries] at (5.5, -0.8) {快速幂迭代步骤（从低位到高位）};

% 表头
\node[font=\scriptsize\bfseries,gray] at (1.2, -1.35) {步骤};
\node[font=\scriptsize\bfseries,gray] at (2.8, -1.35) {当前位};
\node[font=\scriptsize\bfseries,gray] at (4.6, -1.35) {base};
\node[font=\scriptsize\bfseries,gray] at (7.0, -1.35) {操作};
\node[font=\scriptsize\bfseries,gray] at (9.8, -1.35) {result};

\draw[gray!40] (0.5, -1.55) -- (11.5, -1.55);

% 初始
\node[font=\scriptsize] at (1.2, -2.1) {初始};
\node[font=\scriptsize] at (2.8, -2.1) {—};
\node[font=\scriptsize,draw=blue!50,fill=blue!8,rounded corners=1pt,inner sep=2pt]
  at (4.6, -2.1) {base = x};
\node[font=\scriptsize] at (7.0, -2.1) {n = 13 = 1101};
\node[font=\scriptsize,draw=gray!50,fill=white,rounded corners=1pt,inner sep=2pt]
  at (9.8, -2.1) {result = 1};

% 步骤1：bit0=1，result *= x^1
\node[font=\scriptsize] at (1.2, -2.85) {i=0};
\node[font=\scriptsize,draw=blue!60,fill=blue!15,rounded corners=1pt,inner sep=2pt]
  at (2.8, -2.85) {bit=1};
\node[font=\scriptsize] at (4.6, -2.85) {$x^1$};
\node[font=\scriptsize] at (7.0, -2.85) {result $\times$= base ($x^1$)};
\node[font=\scriptsize,draw=orange!70,fill=orange!10,rounded corners=1pt,inner sep=2pt]
  at (9.8, -2.85) {result = $x^1$};

% 步骤2：bit1=0，跳过
\node[font=\scriptsize] at (1.2, -3.6) {i=1};
\node[font=\scriptsize,gray,draw=gray!40,rounded corners=1pt,inner sep=2pt]
  at (2.8, -3.6) {bit=0};
\node[font=\scriptsize] at (4.6, -3.6) {$x^2$};
\node[font=\scriptsize,gray] at (7.0, -3.6) {跳过（bit=0）};
\node[font=\scriptsize,gray] at (9.8, -3.6) {result = $x^1$};

% 步骤3：bit2=1，result *= x^4
\node[font=\scriptsize] at (1.2, -4.35) {i=2};
\node[font=\scriptsize,draw=blue!60,fill=blue!15,rounded corners=1pt,inner sep=2pt]
  at (2.8, -4.35) {bit=1};
\node[font=\scriptsize] at (4.6, -4.35) {$x^4$};
\node[font=\scriptsize] at (7.0, -4.35) {result $\times$= base ($x^4$)};
\node[font=\scriptsize,draw=orange!70,fill=orange!10,rounded corners=1pt,inner sep=2pt]
  at (9.8, -4.35) {result = $x^5$};

% 步骤4：bit3=1，result *= x^8
\node[font=\scriptsize] at (1.2, -5.1) {i=3};
\node[font=\scriptsize,draw=blue!60,fill=blue!15,rounded corners=1pt,inner sep=2pt]
  at (2.8, -5.1) {bit=1};
\node[font=\scriptsize] at (4.6, -5.1) {$x^8$};
\node[font=\scriptsize] at (7.0, -5.1) {result $\times$= base ($x^8$)};
\node[font=\scriptsize,draw=green!60!black,fill=green!10,rounded corners=1pt,inner sep=2pt]
  at (9.8, -5.1) {result = $x^{13}$};

\draw[gray!40] (0.5, -5.45) -- (11.5, -5.45);

%% 分解式说明
\node[draw=gray!40,fill=white,thin,rounded corners=2pt,
      font=\small,inner sep=8pt,align=center]
  at (5.5, -6.5)
  {$x^{13} = x^8 \cdot x^4 \cdot x^1$\quad（对应二进制 \textbf{1}~\textbf{1}~0~\textbf{1} 中值为1的位）\\[4pt]
   base 每步平方：$x \to x^2 \to x^4 \to x^8 \to \ldots$};

%% 流程图（简化树形展示）
\node[font=\small\bfseries] at (5.5, -7.8) {分治递归树（等价视角）};

\node[draw=blue!50,fill=blue!8,rounded corners=2pt,font=\scriptsize,inner sep=4pt]
  (r13) at (5.5, -8.5) {$x^{13}$};
\node[draw=blue!50,fill=blue!8,rounded corners=2pt,font=\scriptsize,inner sep=4pt]
  (r6) at (3.5, -9.5) {$x^6$};
\node[draw=blue!50,fill=blue!8,rounded corners=2pt,font=\scriptsize,inner sep=4pt]
  (r1a) at (7.5, -9.5) {$x^1$（余）};
\node[draw=blue!50,fill=blue!8,rounded corners=2pt,font=\scriptsize,inner sep=4pt]
  (r3) at (2.5, -10.5) {$x^3$};
\node[draw=blue!50,fill=blue!8,rounded corners=2pt,font=\scriptsize,inner sep=4pt]
  (r0) at (4.5, -10.5) {$x^0$（余）};
\node[draw=blue!50,fill=blue!8,rounded corners=2pt,font=\scriptsize,inner sep=4pt]
  (r1b) at (1.8, -11.5) {$x^1$};
\node[draw=green!60!black,fill=green!10,rounded corners=2pt,font=\scriptsize,inner sep=4pt]
  (r1c) at (3.2, -11.5) {$x^1$（余）};

\draw[ptr,gray!60] (r13) -- (r6)  node[midway,font=\tiny,fill=white,inner sep=1pt] {$\div2$};
\draw[ptr,gray!60] (r13) -- (r1a) node[midway,font=\tiny,fill=white,inner sep=1pt] {奇};
\draw[ptr,gray!60] (r6) -- (r3)   node[midway,font=\tiny,fill=white,inner sep=1pt] {$\div2$};
\draw[ptr,gray!60] (r6) -- (r0)   node[midway,font=\tiny,fill=white,inner sep=1pt] {偶};
\draw[ptr,gray!60] (r3) -- (r1b)  node[midway,font=\tiny,fill=white,inner sep=1pt] {$\div2$};
\draw[ptr,gray!60] (r3) -- (r1c)  node[midway,font=\tiny,fill=white,inner sep=1pt] {奇};

%% 算法说明
\node[draw=gray!50,fill=white,rounded corners=3pt,font=\small,
      inner sep=8pt,align=left]
  at (5.5, -13.0)
  {\textbf{快速幂 O(log n) 时间：}\\[4pt]
   迭代法：base=x，每轮 base=base*base；n 为奇时 result*=base；n>>=1\\[2pt]
   递归法：pow(x,n)=pow(x*x,n/2)（n偶），或 x*pow(x,n-1)（n奇）\\[2pt]
   核心：利用二进制分解，$x^n = \prod_{bit_i=1} x^{2^i}$};

\end{tikzpicture}
\end{document}
